% Skenario uji tujuan T-4, pecahan uji_penerimaan
{
\footnotesize
\setlength{\tabcolsep}{3pt}
\Needspace{10\baselineskip}
\begin{longtable}{|>{\centering\arraybackslash}m{1.4cm}|>{\centering\arraybackslash}m{1.75cm}|m{5.9cm}|m{4cm}|}
\caption{Skenario Uji Penerimaan Tujuan T-4}
\label{tab:uji_penerimaan_t4} \\
\hline
\textbf{Skenario} & \textbf{Kebutuhan} & \multicolumn{1}{>{\centering\arraybackslash}m{5.9cm}|}{\textbf{Prosedur Uji}} & \multicolumn{1}{>{\centering\arraybackslash}m{4cm}|}{\textbf{Kriteria Penerimaan}} \\
\hline
\endfirsthead
\hline
\textbf{Skenario} & \textbf{Kebutuhan} & \multicolumn{1}{>{\centering\arraybackslash}m{5.9cm}|}{\textbf{Prosedur Uji}} & \multicolumn{1}{>{\centering\arraybackslash}m{4cm}|}{\textbf{Kriteria Penerimaan}} \\
\hline
\endhead
\endfoot
\hline
\endlastfoot

SK-F-02 & KF-02 & Memeriksa \textit{online feature serving} pada Valkey mengembalikan nilai non-\textit{null} melalui \texttt{online-serving-check.py}. & \textit{Online serving} mengembalikan nilai non-\textit{null} dan kesetaraan nilai \textit{online} dan \textit{offline} terpenuhi secara struktural melalui materialisasi. \\
\hline
SK-F-04 & KF-04 & Menjalankan beban pencarian vektor pada Qdrant melalui \texttt{load/vector-search-latency.js} dengan k6. & Latensi pencarian vektor berada pada orde sub-detik. Pengujian \textit{recall} dilakukan saat koleksi vektor terisi. \\
\hline
SK-F-11 & KF-11 & Membandingkan fitur hasil \textit{batch} (dbt) dan hasil \textit{stream} (Flink) untuk indikator yang sama. & Nilai fitur \textit{batch} dan \textit{stream} setara dalam toleransi yang ditetapkan. \\
\hline
SK-N-01 & KNF-01 & Menjalankan \texttt{load/feast-online-latency.js} dan beban pencarian vektor dengan k6 sambil memantau SLO Sloth. & Latensi p99 berada pada orde sub-detik pada \textit{feature serving} dan \textit{vector search}. \\
\hline
SK-N-02 & KNF-02 & Menjalankan \texttt{load/feast-throughput.js} untuk \textit{serving} dan \texttt{kafka-producer-perf-test} untuk \textit{stream}, sambil memantau \textit{lag}. & \textit{Throughput} mencapai target dan \textit{lag} konsumen tetap terkendali pada beban puncak. \\
\hline
\end{longtable}
}
